\section{Lambda calculus from scratch}
\label{sec:refresher-lambda}

Turing machines make state and local time explicit.  Lambda calculus makes
substitution and composition explicit.  We start with its untyped core,
because fixed-point combinators live there naturally, and only later return to
the typed, fuelled language used by the project.

\subsection{Terms, scope, and beta reduction}

A lambda term is generated by
\[
 t ::= x \mid (\lambda x.t) \mid (t\;t).
\]
A variable occurrence is \emph{bound} if it lies inside a matching
\(\lambda x\); otherwise it is \emph{free}.  For example, in
\(\lambda x.(x\;y)\), the occurrence of \(x\) is bound and \(y\) is free.
Changing a bound name consistently is \emph{alpha-renaming}:
\(\lambda x.x=_{\alpha}\lambda z.z\).  It changes typography, not the term's
binding structure.

Application associates to the left and abstraction extends to the right, so
\(f\;a\;b\) means \((f\;a)\;b\), while \(\lambda x.f\;x\) means
\(\lambda x.(f\;x)\).  The computational rule is beta reduction,
\[
 (\lambda x.t)\;u\longrightarrow_\beta t[x:=u],
\]
where substitution replaces the free occurrences of \(x\) in \(t\) and first
alpha-renames binders when necessary to avoid capture.

\paragraph{Why this matters here.}
Hardwiring a machine input and applying a function are both forms of
substitution.  Our later \(\Specialize\) operation is the source-code version:
it fills typed holes while preserving scope.  Being precise about bound and
free variables is what prevents a verifier description from changing meaning
when code is inserted under a binder.

\paragraph{Worked reduction 1.}
The identity combinator \(I=\lambda x.x\) satisfies
\[
 (\lambda x.x)\;((\lambda z.z)\;w)
 \longrightarrow_\beta (\lambda z.z)\;w
 \longrightarrow_\beta w.
\]

\paragraph{Worked reduction 2: avoiding capture.}
Naively substituting \(y\) for \(x\) in
\((\lambda x.\lambda y.x)\;y\) would incorrectly bind the free \(y\).
Rename the inner binder first:
\[
 (\lambda x.\lambda y.x)\;y
 =_\alpha(\lambda x.\lambda z.x)\;y
 \longrightarrow_\beta\lambda z.y.
\]
The surviving \(y\) is free, as it must be.

A \emph{beta-redex} is a subterm of the form \((\lambda x.t)u\).  A
\emph{normal form} contains no beta-redex.  Some terms have none: with
\(\Omega=(\lambda x.xx)(\lambda x.xx)\), one has
\(\Omega\to_\beta\Omega\) forever.  The Church--Rosser, or confluence,
theorem states that if one term reduces by different paths to \(u\) and \(v\),
then \(u\) and \(v\) have a common reduct.  In particular, a normal form, if
it exists, is unique up to alpha-renaming.  We use this theorem as orientation
and do not prove it here.

\paragraph{Why this matters here.}
Confluence says that the mathematical result does not depend on which redex
is chosen, provided reduction reaches a normal form.  Resource accounting is
stricter: different strategies may take very different numbers of steps, and
one may diverge while another reaches the answer.  The project therefore
fixes an evaluator instead of treating all beta paths as operationally equal.

\paragraph{Worked strategy example.}
For \(K=\lambda x.\lambda y.x\),
\[
 K\;I\;\Omega\longrightarrow_\beta I
\]
under normal order, which reduces the leftmost outermost redex first and never
evaluates the unused argument.  Call-by-value first tries to turn \(\Omega\)
into a value and therefore loops.  This difference is exactly why the usual
\(Y\) combinator needs an eta-delayed cousin under call-by-value.

\subsection{Church data: behavior as representation}

Pure lambda calculus has no primitive bits, tuples, or integers.  Church
encodings represent data by what eliminator it applies.  These encodings are
not the project's efficient runtime representation; they show that the bare
calculus is computationally expressive.

\paragraph{Booleans.}
Define
\[
 \mathsf{true}=\lambda t.\lambda f.t,
 \qquad \mathsf{false}=\lambda t.\lambda f.f.
\]
They choose the first or second continuation.  For example,
\(\mathsf{true}\;A\;B\to_\beta A\), whereas
\(\mathsf{false}\;A\;B\to_\beta B\).

\paragraph{Why booleans matter here.}
A verifier ultimately returns one acceptance bit.  In the resource language
we use a primitive bit sort, but its conditional has precisely the selection
behavior exhibited by these terms.

\paragraph{If--then--else.}
Define
\(\mathsf{if}=\lambda b.\lambda t.\lambda e.b\;t\;e\).  Then
\[
 \mathsf{if}\;\mathsf{false}\;A\;B
 \to_\beta\mathsf{false}\;A\;B
 \to_\beta B.
\]

\paragraph{Why the conditional matters here.}
Every decider is built from branches on parsed bits, timeout tests, and local
predicates.  Fixing their evaluation order makes ``the branch not taken''
operationally meaningful, especially when it contains a recursive call.

\paragraph{Pairs.}
Let
\[
 \mathsf{pair}=\lambda a.\lambda b.\lambda p.p\;a\;b,
 \quad \mathsf{fst}=\lambda p.p\;\mathsf{true},
 \quad \mathsf{snd}=\lambda p.p\;\mathsf{false}.
\]
Writing \(\langle A,B\rangle=\mathsf{pair}\;A\;B\),
\[
 \mathsf{fst}\;\langle A,B\rangle
 \to_\beta\mathsf{true}\;A\;B\to_\beta A.
\]

\paragraph{Why pairs matter here.}
Machine configurations, verifier descriptions, and sampler/decider outputs
are finite products.  The project gives tuples primitive storage so that
projection has an explicit small cost, but the lambda interface is the same.

\paragraph{Numerals.}
The Church numeral \(n\) is
\[
 \overline n=\lambda f.\lambda x.f^n x,
 \quad
 \overline0=\lambda f.\lambda x.x,
 \quad
 \overline2=\lambda f.\lambda x.f(fx).
\]
It represents iteration: \(\overline2\;s\;z\to_\beta s(sz)\).

\paragraph{Why numerals matter here.}
An index \(n\), a timeout, and a fuel counter are all iterated control data.
Church numerals prove representability, while binary primitive naturals in
Part~II prevent a unary blow-up in the resource bounds.

\paragraph{Successor.}
Define
\(\mathsf{succ}=\lambda n.\lambda f.\lambda x.f(nfx)\).  Then
\[
 \mathsf{succ}\;\overline2
 \to_\beta\lambda f.\lambda x.f(\overline2 f x)
 \to_\beta\lambda f.\lambda x.f(f(fx))=\overline3.
\]

\paragraph{Why successor matters here.}
It is the simplest example of a computable map on an encoding.  Later
compilers similarly build a new syntax tree around an old one and prove a
size increase, rather than treating the new program as an abstract function.

\paragraph{Addition.}
Define
\(\mathsf{add}=\lambda m.\lambda n.\lambda f.\lambda x.mf(nfx)\).  Thus
\[
 \mathsf{add}\;\overline1\;\overline2\;f\;x
 \to_\beta \overline1 f(\overline2 f x)
 \to_\beta f(f(fx)),
\]
so the result is \(\overline3\).

\paragraph{Why addition matters here.}
Fuel bounds and encoded addresses are assembled from smaller bounds.  This
example reminds us that arithmetic is itself computation; Part~II uses
efficient binary primitives and charges for their reference machines.

\paragraph{Predecessor by the pair trick.}
Predecessor is subtler because a Church numeral only says ``iterate.''  Let
\[
 \Phi=\lambda p.\mathsf{pair}\;(\mathsf{snd}\;p)
                 \;(\mathsf{succ}(\mathsf{snd}\;p)),
 \qquad
 \mathsf{pred}=\lambda n.\mathsf{fst}
   (n\;\Phi\;(\mathsf{pair}\;\overline0\;\overline0)).
\]
Each \(\Phi\)-step maps \(\langle j-1,j\rangle\) to
\(\langle j,j+1\rangle\).  The complete small calculation for three is
\[
 \langle0,0\rangle\xrightarrow{\Phi}\langle0,1\rangle
 \xrightarrow{\Phi}\langle1,2\rangle
 \xrightarrow{\Phi}\langle2,3\rangle
 \xrightarrow{\mathsf{fst}}2.
\]
Hence \(\mathsf{pred}\;\overline3\to_\beta^*\overline2\).

\paragraph{Why predecessor matters here.}
A timeout counts down, so total bounded simulation needs a predecessor and a
zero test even when the simulated program diverges.  The pair construction is
a miniature instance of adding auxiliary state to make a desired update
locally computable.

\paragraph{Lists.}
Define a fold encoding
\[
 \mathsf{nil}=\lambda c.\lambda z.z,
 \qquad
 \mathsf{cons}=\lambda h.\lambda t.\lambda c.\lambda z.
                    c\;h\;(t\;c\;z).
\]
Then the list \([A,B]=\mathsf{cons}\;A\;(\mathsf{cons}\;B\;\mathsf{nil})\)
satisfies
\[
 [A,B]\;C\;Z\longrightarrow_\beta^* C\;A\;(C\;B\;Z).
\]

\paragraph{Why lists matter here.}
Tapes, syntax children, environments, and continuation stacks are all finite
sequences with potentially unbounded length.  A functional list supplies the
shape; the resource IR adds heap pointers and charged operations so that the
cost of traversing it is visible.

\subsection{Recursion from self-application}

Suppose \(F\) maps a proposed recursive procedure to its body.  We want a term
\(f\) satisfying
\[
 f=Ff.
\]
Think of the occurrence of \(f\) on the right as a hole.  If the desired term
can be written as self-application \(h\;h\), then the equation asks for
\(h\;h=F(h\;h)\).  Choose
\[
 h=\lambda x.F(x\;x).
\]
Substituting this \(h\) into \(h\;h\) yields the fixed-point combinator
\[
 Y=\lambda F.(\lambda x.F(x\;x))(\lambda x.F(x\;x)).
\]
The promised equation is a direct beta calculation:
\[
\begin{aligned}
 Y\;F
 &\longrightarrow_\beta
 (\lambda x.F(xx))(\lambda x.F(xx))\\
 &\longrightarrow_\beta
 F\bigl((\lambda x.F(xx))(\lambda x.F(xx))\bigr)\\
 &=F(YF).
\end{aligned}
\]
There is no recursive name in the finite syntax of \(Y\); recursion appears
when two copies are applied.

\paragraph{Why this matters here.}
The lambda fixed point is the term-level analogue of Kleene's theorem.  Both
build a finite object whose unfolding supplies that object's own behavior.
Our IR will make the analogous code link explicit and constant-size, because
blindly expanding \(YF\) would obscure description length and fuel.

\paragraph{Worked fixed point.}
Take \(F=\lambda r.\lambda n.\mathsf{if}\;(\mathsf{iszero}\;n)
\;\overline1\;(\mathsf{mul}\;n\;(r(\mathsf{pred}\;n)))\).  Under normal-order
reduction, \(YF\;\overline3\) unfolds only when the selected branch needs it:
\[
\begin{aligned}
 YF\;3&\to F(YF)\;3\\
 &\to 3\cdot(YF\;2)
 \to 3\cdot2\cdot(YF\;1)\\
 &\to3\cdot2\cdot1\cdot(YF\;0)
 \to3\cdot2\cdot1\cdot1=6.
\end{aligned}
\]
Here the arithmetic names abbreviate their Church encodings; every displayed
arrow stands for their finite beta reductions.

Call-by-value evaluates an argument before entering the function body, so
\(YF\) tries to expand itself forever.  Eta-delay the recursive occurrence:
\[
 Z=\lambda F.
 (\lambda x.F(\lambda u.xxu))
 (\lambda x.F(\lambda u.xxu)).
\]
Now
\[
 ZF\to_\beta F(\lambda u.(ZF)u),
\]
and the recursive computation remains under a lambda until an actual argument
arrives.  Under call-by-value the conditional branches must likewise be
thunked.  With
\[
 F_v=\lambda r.\lambda n.
 \bigl(\mathsf{if}\;(\mathsf{iszero}\;n)\;
   (\lambda u.\overline1)\;
   (\lambda u.\mathsf{mul}\;n\;(r(\mathsf{pred}\;n)))\bigr)\;I,
\]
only the selected branch is applied to the dummy value \(I\), and
\(ZF_v\;\overline3\to_\beta^*\overline6\) under call-by-value.

\paragraph{Why evaluation strategy matters here.}
The project evaluator is deterministic call-by-value, so its fixed-point
constructor must behave like \(Z\), not raw \(Y\).  Once the strategy is
fixed, ``number of reductions'' becomes an honest resource that can be
translated to Turing-machine time.

\subsection{De Bruijn indices and quotation}

Names are pleasant for handwriting but awkward for a compiler: alpha-equal
terms have many spellings, and substitution must keep inventing fresh names.
De Bruijn notation replaces a bound occurrence by the number of binders
between it and its binder, counting the nearest binder as zero.  Thus
\[
 \lambda x.\lambda y.x\;y
 \quad\text{becomes}\quad
 \lambda.\lambda.(1\;0).
\]
Both \(\lambda a.\lambda b.a\;b\) and
\(\lambda x.\lambda y.x\;y\) serialize to the same tree.

\paragraph{Why this matters here.}
The IR needs canonical bytes and capture-free specialization.  De Bruijn
indices remove alpha-renaming from the representation: a variable is a
lexical address, and inserting a closed term cannot capture one of its free
names because it has none.

\paragraph{Worked index lookup.}
In \(\lambda.\lambda.(1\;0)\), applying the outer function to \(A\) and the
inner function to \(B\) creates environment frames \([B],[A]\).  Index zero
selects \(B\); index one crosses one frame and selects \(A\).  The body is
therefore \(A\;B\), just as the named term says.

To \emph{quote} a term is to treat its syntax rather than its value as data.
A term is a finite ordered tree: each node has a constructor tag
(variable, lambda, application, and so on), finite payloads, and finitely many
children.  Prefix-free integer and byte encodings serialize that tree to a
unique bit string \(\operatorname{ser}(t)\).  We define
\[
 |t|:=|\operatorname{ser}(t)|.
\]
Quotation does not claim to compute what \(t\) will do.  It merely reifies the
finite tree so an evaluator or compiler can inspect it.

\paragraph{Why this matters here.}
\(\Compress\) must accept code, measure its length, insert static parameters,
and return new code.  A runtime closure with a hidden host-language
environment is not adequate; a quoted, closed tree is.  Quotation is the
lambda-calculus version of writing \(\langle M\rangle\) on a universal
machine's input tape.

\paragraph{Worked quotation.}
The named identity terms \(\lambda x.x\) and \(\lambda z.z\) both become the
de Bruijn tree \(\mathsf{Lambda}(\mathsf{BoundVar}(0))\).  Quoting either
produces the same canonical bytes.  Evaluating those bytes on \(A\) returns
\(A\); copying the bytes returns code, not \(A\).

\section{Church--Turing as an engineering statement}
\label{sec:refresher-engineering}

The Church--Turing thesis is often remembered as a slogan about which
functions are computable.  Here we need a more prosaic engineering reading:
the two concrete programming media can compile into one another, and the
compilers preserve both finite descriptions and feasible bounds up to fixed
polynomials.

\paragraph{Turing machine to lambda term.}
Represent each tape as a functional zipper, represent a configuration as a
tuple, and write a fixed recursive loop which reads the hardwired transition
table and produces the next configuration.  One machine step becomes a
bounded collection of list, tuple, and table operations.  The term contains
one copy of \(\langle M\rangle\), so its serialized size is linear in
\(|M|\); a \(T\)-step machine run becomes a polynomial number of evaluator
steps.  Theorem~\ref{thm:tm-to-l} gives the formal bound.

\paragraph{Why this direction matters here.}
It shows that replacing a paper verifier by an \(\mathcal L\)-term loses none
of its behavior and makes no hidden exponential copy of its transition table.

\paragraph{Worked direction.}
For \(M_=\), the recursive loop starts from the four-tape tuple, looks up the
equal-input row, then the copy row, then the halt row.  Three loop iterations
produce the Church-level bit one, with fixed overhead for tuple and list
operations.

\paragraph{Lambda term to Turing machine.}
Write a fixed evaluator whose work tape stores the parsed syntax tree, an
environment, a continuation stack, and a fuel counter.  Each evaluation rule
is a finite transition routine.  A run using \(F\) charged reductions is
simulated in time polynomial in the code length, input length, and \(F\).
Hardwiring the quoted term into this evaluator gives a machine description
linear in the serialized term.  Theorem~\ref{thm:l-to-tm} makes this precise.

\paragraph{Why this direction matters here.}
The paper's theorems are stated for Turing-machine verifiers.  Compiling a
resource-accounted term back to that model is what licenses using the lambda
IR without silently changing the theorem being applied.

\paragraph{Worked direction.}
For the quoted term \((\lambda x.x)\;1\), the evaluator parses an application
node, builds the identity closure, places the bit in its environment, looks up
index zero, and emits one.  A Turing machine can represent every one of those
finite records on its work tape.

The promised dictionary is therefore:
\begin{center}
\small
\begin{tabularx}{0.94\textwidth}{P{0.29\textwidth}P{0.29\textwidth}X}
\toprule
Turing-machine notion & Lambda/IR notion & Operative correspondence \\
\midrule
Program & Closed term & Finite rule table versus finite syntax tree \\
Input tape contents & Argument term/value & Tuple components keep the inputs
separate \\
Machine transitions and time & Fuel/charged reduction steps & Each simulates
the other with fixed polynomial overhead \\
\(\langle M\rangle\), \(|M|\) & Serialized quoted tree, \(|t|\) & Canonical
finite data and its bit length \\
Universal machine & Evaluator & One fixed program interprets supplied code \\
\(s\)-\(m\)-\(n\) hardwiring & \(\Specialize\) & A computable map from code
and static data to new code \\
Kleene recursion theorem & \(Y\)/\(Z\), or the IR's \(\mathsf{Fix}\) & A
finite behavioral fixed point \\
\bottomrule
\end{tabularx}
\end{center}

\paragraph{What the dictionary does not say.}
It does not assert step-for-step equality, equal constants, or equal
description strings.  It asserts a pair of explicit compilers, semantic
agreement on each input, linear description growth, and polynomial running
time overhead.  Those are exactly the invariants needed when an exponent
\(\lambda\) is enlarged by a universal constant.

\section{Why this machinery appears in a paper about entangled provers}
\label{sec:refresher-why}

A nonlocal game is an analytic object: its value optimizes over strategies,
and in the entangled setting those strategies use states and measurements.
Nevertheless, the \emph{outer construction} of the game is algorithmic.  The
compression theorem starts from the finite description of a normal-form
verifier \(V=(S,D)\), together with a resource exponent \(\lambda\), and
returns another finite verifier description.  Its hypotheses constrain both
\(|V|\) and worst-case running times.  Consequently, extensional statements
such as ``\(D\) decides the same predicate'' are not enough: the proof must
know which code is embedded, how long it is, and how expensive its universal
simulation will be.

At the description level, the main pipeline is
\[
 \langle V\rangle
 \xrightarrow{\ \Introspect\ }
 \langle V^{(1)}\rangle
 \xrightarrow{\ \AnswerReduce\ }
 \langle V^{(2)}\rangle
 \xrightarrow{\ \Repeat\ }
 \langle V^{(3)}\rangle.
\]
Each arrow is a terminating computable map on finite strings.  Introspection
embeds a description so a verifier can check a smaller representation of its
own computation.  Answer reduction constructs a bounded-trace circuit and
then a PCP-style verifier.  Repetition constructs several coordinated copies
and an acceptance rule.  The composite is \(\Compress\).  Polynomial
bookkeeping proves that its output is still a legal bounded verifier.

\paragraph{Why the fixed point enters.}
For a supplied machine \(M\) and level parameter \(L\), the Halting verifier
must arrange that its decider is the one appearing inside the verifier fed to
\(\Compress\).  This is a fixed-point equation on descriptions.  Kleene's
theorem, the explicit \(\mathcal F(\mathcal F,M,L,\ldots)\) program, and the
lambda fixed-point constructor are three presentations of the same outer
operation.  The point is not mystical self-awareness: it is a finite compiler
which prints code containing an appropriate code literal.

\paragraph{Worked pipeline fragment.}
Suppose the input decider has code \(d\).  Answer reduction first constructs a
succinct circuit whose hardwired local transition rule is read from \(d\).
It then emits a new decider code \(d'\) that reconstructs and checks the
corresponding local constraints.  The transformation \(d\mapsto d'\) can be
run by a Turing machine, and its time and output length can be bounded before
either decider is executed.  Compression composes several transformations of
exactly this kind.

\paragraph{What is not computability theory.}
The dictionary stops at the boundary between constructing a game and proving
what quantum strategies can do in it.  The low-degree test's soundness,
rigidity of observables, state-dependent approximation estimates,
oracularization against entangled strategies, entanglement-dimension bounds,
and the quantitative preservation of completeness and soundness are analytic
statements.  They do not follow from universality, quotation, or a fixed-point
theorem.  A lambda implementation can expose the verifier transformation and
check finite algebraic identities; it cannot by that fact alone prove the
operator-algebraic leaves.

\paragraph{Why this separation matters here.}
It prevents two opposite mistakes.  One is to treat source manipulation as
informal and accidentally hide a noncomputable or super-polynomial step.  The
other is to imagine that a clean self-interpreter settles the quantum
soundness argument.  The paper needs both layers: computability theory with
polynomial bookkeeping on the outside, and the genuinely quantum analytic
theorems on the inside.  The resource-accounted description language in
Part~II is designed for the first layer and records exactly where it must call
